\chapter{2010年江苏卷}

\section{填空题}

\begin{problem}
设集合 \(A=\{-1,1,3\}\)，\(B=\{a+2,a^2+4\}\)，\(A\cap B=\{3\}\)，则实数 \(a=\)\fillinblank{}.
\end{problem}

\begin{answer}
1
\end{answer}

\begin{problem}
设复数 \(z\) 满足 \(z(2-3\i)=6+4\i\)（其中 \(\i\) 为虚数单位），则 \(z\) 的模为\fillinblank{}.
\end{problem}

\begin{answer}
2
\end{answer}

\begin{problem}
盒子中有大小相同的 \(3\) 只白球，\(1\) 只黑球，若从中随机地摸出两只球，则它们颜色不同的概率是\fillinblank{}.
\end{problem}

\begin{answer}
\(\frac{1}{2}\)
\end{answer}

\begin{problem}
某棉纺厂为了了解一批棉花的质量，从中随机抽取了 \(100\) 根棉花纤维的长度（棉花纤维的长度是棉花质量的重要指标），所得数据都在区间 \([5,40]\) 中，其频率分布直方图如图所示，则其抽样的 \(100\) 根中，有\fillinblank{} 根棉花纤维的长度小于 \(20 \mathrm{~mm}\).

\bitmapfigure[max width=0.42\linewidth]{img/2010/jiangsu/q4_fig1.png}
\end{problem}

\begin{answer}
30
\end{answer}

\begin{problem}
设函数 \(f(x)=x\left(\e^x+a\e^{-x}\right)\)，\(x\in\R\)，是偶函数，则实数 \(a\) 的值为\fillinblank{}.
\end{problem}

\begin{answer}
\(-1\)
\end{answer}

\begin{problem}
在平面直角坐标系 \(xOy\) 中，已知双曲线 \(\frac{x^2}{4}-\frac{y^2}{12}=1\) 上一点 \(M\) 的横坐标为 \(3\)，则点 \(M\) 到此双曲线的右焦点的距离为\fillinblank{}.
\end{problem}

\begin{answer}
4
\end{answer}

\begin{problem}
如图是一个算法流程图，则输出 \(S\) 的值是\fillinblank{}.

\bitmapfigure[max width=0.34\linewidth]{img/2010/jiangsu/q7_fig1.png}
\end{problem}

\begin{answer}
63
\end{answer}

\begin{problem}
函数 \(y=x^2\)（\(x>0\)）的图象在点 \(\left(a_k,a_k^2\right)\) 处的切线与 \(x\) 轴交点的横坐标为 \(a_{k+1}\)，其中 \(k\in\N^*\). 若 \(a_1=16\)，则 \(a_1+a_3+a_5\) 的值是\fillinblank{}.
\end{problem}

\begin{answer}
21
\end{answer}

\begin{problem}
在平面直角坐标系 \(xOy\) 中，已知圆 \(x^2+y^2=4\) 上有且仅有四个点到直线 \(12x-5y+c=0\) 的距离为 \(1\)，则实数 \(c\) 的取值范围是\fillinblank{}.
\end{problem}

\begin{answer}
\((-13,13)\)
\end{answer}

\begin{problem}
定义在区间 \(\left(0,\frac{\pi}{2}\right)\) 上的函数 \(y=6\cos x\) 的图象与 \(y=5\tan x\) 的图象交于点 \(P\)，过点 \(P\) 作 \(x\) 轴的垂线，垂足为 \(P_1\)，直线 \(PP_1\) 与 \(y=\sin x\) 的图象交于点 \(P_2\)，则线段 \(P_1P_2\) 的长为\fillinblank{}.
\end{problem}

\begin{answer}
\(\frac{2}{3}\)
\end{answer}

\begin{problem}
已知函数 \(f(x)=\begin{cases}
x^2+1,&x\geqslant 0,\\
1,&x<0,
\end{cases}\) 则满足不等式 \(f(1-x^2)>f(2x)\) 的 \(x\) 的取值范围是\fillinblank{}.
\end{problem}

\begin{answer}
\((-1,\sqrt{2}-1)\)
\end{answer}

\begin{problem}
设实数 \(x\)，\(y\) 满足 \(3\leqslant xy^2\leqslant 8\)，\(4\leqslant \frac{x^2}{y}\leqslant 9\)，则 \(\frac{x^3}{y^4}\) 的最大值是\fillinblank{}.
\end{problem}

\begin{answer}
27
\end{answer}

\begin{problem}
在锐角 \(\triangle ABC\) 中，角 \(A\)，\(B\)，\(C\) 的对边分别为 \(a\)，\(b\)，\(c\)，若 \(\frac{b}{a}+\frac{a}{b}=6\cos C\)，则 \(\frac{\tan C}{\tan A}+\frac{\tan C}{\tan B}\) 的值是\fillinblank{}.
\end{problem}

\begin{answer}
4
\end{answer}

\begin{problem}
将边长为 \(1 \mathrm{~m}\) 的正三角形薄铁皮沿一条平行于某边的直线剪成两块，其中一块是梯形，记 \(s=\dfrac{\left(\text{梯形的周长}\right)^2}{\text{梯形的面积}}\)，则 \(s\) 的最小值是\fillinblank{}.
\end{problem}

\begin{answer}
\(\frac{32\sqrt{3}}{3}\)
\end{answer}

\section{解答题}

\begin{problem}
在平面直角坐标系 \(xOy\) 中，已知点 \(A(-1,-2)\)，\(B(2,3)\)，\(C(-2,-1)\).

\begin{enumerate}
\item 求以线段 \(AB\)，\(AC\) 为邻边的平行四边形的两条对角线的长；
\item 设实数 \(t\) 满足 \(\left(\overrightarrow{AB}-t\overrightarrow{OC}\right)\cdot \overrightarrow{OC}=0\)，求 \(t\) 的值.
\end{enumerate}
\end{problem}

\begin{solution}
\begin{enumerate}
\item 由题意
\[
\overrightarrow{AB}=(3,5),\qquad \overrightarrow{AC}=(-1,1).
\]
以 \(AB\)，\(AC\) 为邻边的平行四边形的两条对角线所对应的向量分别为
\[
\overrightarrow{AB}+\overrightarrow{AC}=(2,6),\qquad \overrightarrow{AB}-\overrightarrow{AC}=(4,4).
\]
因此两条对角线的长分别为
\[
\sqrt{2^2+6^2}=2\sqrt{10},\qquad \sqrt{4^2+4^2}=4\sqrt{2}.
\]

\item 由 \(O(0,0)\)，\(C(-2,-1)\) 得
\[
\overrightarrow{OC}=(-2,-1),\qquad \overrightarrow{AB}=(3,5).
\]
代入已知条件，得
\[
\left(\overrightarrow{AB}-t\overrightarrow{OC}\right)\cdot\overrightarrow{OC}
=\overrightarrow{AB}\cdot\overrightarrow{OC}-t\left|\overrightarrow{OC}\right|^2
=-11-5t=0.
\]
所以 \(t=-\dfrac{11}{5}\).
\end{enumerate}
\end{solution}

\begin{problem}
如图，在四棱锥 \(P-ABCD\) 中，\(PD \perp\) 平面 \(ABCD\)，\(PD=DC=BC=1\)，\(AB=2\)，\(AB \parallel DC\)，\(\angle BCD=90^\circ\).

\begin{enumerate}
\item 求证：\(PC \perp BC\)；
\item 求点 \(A\) 到平面 \(PBC\) 的距离.
\end{enumerate}

\bitmapfigure[max width=0.42\linewidth]{img/2010/jiangsu/q16_fig1.png}
\end{problem}

\begin{solution}
\begin{enumerate}
\item 因为 \(PD\perp\) 平面 \(ABCD\)，所以 \(PD\perp BC\). 又由 \(\angle BCD=90^\circ\)，得 \(DC\perp BC\). 直线 \(PD\)，\(DC\) 是平面 \(PDC\) 内相交的两条直线，所以 \(BC\perp\) 平面 \(PDC\). 因为 \(PC\subset\) 平面 \(PDC\)，故 \(PC\perp BC\).

\item 建立如图所示的直角坐标系，取
\[
D(0,0,0),\qquad C(1,0,0),\qquad B(1,1,0),\qquad P(0,0,1).
\]
由 \(AB\parallel DC\)，\(AB=2\)，可取 \(A(-1,1,0)\). 于是平面 \(PBC\) 的一个法向量为
\[
\overrightarrow{PB}\times\overrightarrow{PC}
=(1,1,-1)\times(1,0,-1)=(-1,0,-1).
\]
因此平面 \(PBC\) 的方程为
\[
x+z-1=0.
\]
点 \(A(-1,1,0)\) 到该平面的距离为
\[
\frac{|-1+0-1|}{\sqrt{1^2+1^2}}=\sqrt{2}.
\]
所以点 \(A\) 到平面 \(PBC\) 的距离为 \(\sqrt{2}\).
\end{enumerate}
\end{solution}

\begin{problem}
某兴趣小组测量电视塔 \(AE\) 的高度 \(H\)（单位：\(\mathrm{m}\)）. 如示意图，垂直放置的标杆 \(BC\) 的高度 \(h=4 \mathrm{~m}\)，仰角 \(\angle ABE=\alpha\)，\(\angle ADE=\beta\).

\begin{enumerate}
\item 该小组已经测得一组 \(\alpha\)，\(\beta\) 的值，算出了 \(\tan\alpha=1.24\)，\(\tan\beta=1.20\)，请据此算出 \(H\) 的值；
\item 该小组分析若干测得的数据后，发现适当调整标杆到电视塔的距离 \(d\)（单位：\(\mathrm{m}\)），使 \(\alpha\) 与 \(\beta\) 之差较大，可以提高测量精度. 若电视塔实际高度为 \(125 \mathrm{~m}\)，试问 \(d\) 为多少时，\(\alpha-\beta\) 最大.
\end{enumerate}

\bitmapfigure[max width=0.42\linewidth]{img/2010/jiangsu/q17_fig1.png}
\end{problem}

\begin{solution}
由图，点 \(D\)，\(B\)，\(A\) 共线，点 \(D\)，\(C\)，\(E\) 共线，且 \(BC\perp DA\)，\(AE\perp DA\). 记 \(BD=x\)，则 \(x>0\)，\(BA=d\)，\(BC=h\)，\(AE=H\).

\begin{enumerate}
\item 在直角三角形 \(DBC\) 和 \(DAE\) 中，有 \(\tan\beta=\frac{h}{x}=\frac{H}{x+d}\)，且 \(\tan\alpha=\frac{H}{d}\).
由前一个等式得 \(h=\tan\beta\,x\)，\(H=\tan\beta(x+d)\)，两式相减得 \(H-h=\tan\beta\,d\). 又因为 \(d=H/\tan\alpha\)，所以
\(H-h=\frac{\tan\beta}{\tan\alpha}H\)，从而 \(H=\frac{h\tan\alpha}{\tan\alpha-\tan\beta}\).
代入 \(h=4\)，\(\tan\alpha=1.24\)，\(\tan\beta=1.20\)，得
\[
H=\frac{4\times1.24}{1.24-1.20}=124\ \mathrm{m}.
\]

\item 因为直角三角形 \(DBC\) 和 \(DAE\) 相似，所以
\[
\frac{BD}{BA+BD}=\frac{BC}{AE}=\frac{h}{H}.
\]
解得
\[
BD=\frac{hd}{H-h}.
\]
于是 \(\tan\alpha=\frac{H}{d}=\frac{125}{d}\)，\(\tan\beta=\frac{h}{BD}=\frac{H-h}{d}=\frac{121}{d}\).
由于 \(d>0\)，且 \(\alpha\)，\(\beta\) 都是锐角，令
\[
\varphi(d)=\alpha-\beta=\arctan\frac{125}{d}-\arctan\frac{121}{d}.
\]
则
\[
\varphi'(d)=-\frac{125}{d^2+125^2}+\frac{121}{d^2+121^2}
=\frac{4(125\times121-d^2)}{(d^2+125^2)(d^2+121^2)}.
\]
因此，当 \(0<d<\sqrt{125\times121}\) 时，\(\varphi'(d)>0\)；当 \(d>\sqrt{125\times121}\) 时，\(\varphi'(d)<0\). 所以 \(\varphi(d)\) 在
\[
d=\sqrt{125\times121}=5\sqrt{605}\ \mathrm{m}
\]
时取得最大值.
\end{enumerate}
\end{solution}

\begin{problem}
在平面直角坐标系 \(xOy\) 中，如图，已知椭圆 \(\frac{x^2}{9}+\frac{y^2}{5}=1\) 的左、右顶点为 \(A\)，\(B\)，右焦点为 \(F\). 过点 \(T(t,m)\) 的直线 \(TA\)，\(TB\) 与此椭圆分别交于点 \(M(x_1,y_1)\)，\(N(x_2,y_2)\)，其中 \(m>0\)，\(y_1>0\)，\(y_2<0\).

\begin{enumerate}
\item 设动点 \(P\) 满足 \(PF^2-PB^2=4\)，求点 \(P\) 的轨迹；
\item 设 \(x_1=2\)，\(x_2=\frac{1}{3}\)，求点 \(T\) 的坐标；
\item 设 \(t=9\)，求证：直线 \(MN\) 必过 \(x\) 轴上的一定点（其坐标与 \(m\) 无关）.
\end{enumerate}

\bitmapfigure[max width=0.42\linewidth]{img/2010/jiangsu/q18_fig1.png}
\end{problem}

\begin{solution}
\begin{enumerate}
\item 由椭圆方程可知 \(a=3\)，\(b=\sqrt{5}\)，所以右焦点为 \(F(2,0)\)，且 \(B(3,0)\). 设 \(P(x,y)\)，则
\[
PF^2-PB^2=\left((x-2)^2+y^2\right)-\left((x-3)^2+y^2\right)=2x-5.
\]
由题设 \(2x-5=4\)，得 \(x=\dfrac{9}{2}\). 因而点 \(P\) 的轨迹为直线 \(x=\dfrac{9}{2}\).

\item 将 \(x_1=2\) 代入椭圆方程，且由 \(y_1>0\)，得
\[
y_1=\frac{5}{3}.
\]
又 \(A(-3,0)\)，所以直线 \(AM\) 的方程为
\[
y=\frac{x+3}{3}.
\]
将 \(x_2=\dfrac{1}{3}\) 代入椭圆方程，且由 \(y_2<0\)，得
\[
y_2=-\frac{20}{9}.
\]
又 \(B(3,0)\)，所以直线 \(BN\) 的方程为
\[
y=\frac{5}{6}(x-3).
\]
联立两直线方程，得
\[
\frac{x+3}{3}=\frac{5}{6}(x-3),\qquad x=7,\qquad y=\frac{10}{3}.
\]
所以 \(T\) 的坐标为 \(\left(7,\dfrac{10}{3}\right)\).

\item 当 \(t=9\) 时，直线 \(AT\)，\(BT\) 的方程分别为
\[
y=\frac{m}{12}(x+3),\qquad y=\frac{m}{6}(x-3).
\]
设 \(M(x_M,y_M)\). 将第一条直线代入椭圆方程，得到
\[
\frac{x^2}{9}+\frac{m^2(x+3)^2}{720}=1.
\]
其中一个根为 \(x=-3\)，另一个根满足
\[
x_M+3=\frac{480}{m^2+80}.
\]
所以
\[
x_M=\frac{3(80-m^2)}{m^2+80},\qquad y_M=\frac{40m}{m^2+80}.
\]

同理，设 \(N(x_N,y_N)\)，将第二条直线代入椭圆方程，其中一个根为 \(x=3\)，另一个根满足
\[
x_N-3=-\frac{120}{m^2+20}.
\]
所以
\[
x_N=\frac{3(m^2-20)}{m^2+20},\qquad y_N=-\frac{20m}{m^2+20}.
\]
由于 \(m>0\)，有 \(y_M>0\)，\(y_N<0\)，与题设相符.

设直线 \(MN\) 与 \(x\) 轴交于 \(Q(x_0,0)\). 由两点式，
\[
x_0=\frac{x_My_N-x_Ny_M}{y_N-y_M}.
\]
代入上式中的坐标，分子为
\[
\frac{-60m(m^2+40)}{(m^2+80)(m^2+20)},
\]
分母为
\[
\frac{-60m(m^2+40)}{(m^2+80)(m^2+20)}.
\]
因此 \(x_0=1\)，即直线 \(MN\) 恒过定点 \(Q(1,0)\)，该点坐标与 \(m\) 无关.
\end{enumerate}
\end{solution}

\begin{problem}
设各项均为正数的数列 \(\{a_n\}\) 的前 \(n\) 项和为 \(S_n\). 已知 \(2a_2=a_1+a_3\)，数列 \(\left\{\sqrt{S_n}\right\}\) 是公差为 \(d\) 的等差数列.

\begin{enumerate}
\item 求数列 \(\{a_n\}\) 的通项公式（用 \(n\)，\(d\) 表示）；
\item 设 \(c\) 为实数，对满足 \(m+n=3k\) 且 \(m\neq n\) 的任意正整数 \(m\)，\(n\)，\(k\)，不等式 \(S_m+S_n>cS_k\) 都成立. 求证：\(c\) 的最大值为 \(\frac{9}{2}\).
\end{enumerate}
\end{problem}

\begin{solution}
\begin{enumerate}
\item 记 \(\sqrt{S_1}=r\). 因为各项 \(a_n\) 均为正数，所以 \(S_n\) 严格递增，从而等差数列 \(\{\sqrt{S_n}\}\) 的公差满足 \(d>0\)，且
\[
\sqrt{S_n}=r+(n-1)d.
\]
由 \(a_1=S_1\)，\(a_2=S_2-S_1\)，\(a_3=S_3-S_2\) 及 \(2a_2=a_1+a_3\)，得
\[
3S_2=3S_1+S_3.
\]
代入 \(S_1=r^2\)，\(S_2=(r+d)^2\)，\(S_3=(r+2d)^2\)，得
\[
(r+2d)^2-3(r+d)^2+3r^2=(r-d)^2=0.
\]
所以 \(r=d\)，即
\[
S_n=n^2d^2.
\]
于是 \(a_1=S_1=d^2\)，且当 \(n\geqslant2\) 时
\[
a_n=S_n-S_{n-1}=n^2d^2-(n-1)^2d^2=(2n-1)d^2.
\]
所以对所有正整数 \(n\)，都有 \(a_n=(2n-1)d^2\).

\item 由 \(m+n=3k\)，有
\[
S_m+S_n=d^2(m^2+n^2)=\frac{d^2}{2}\left((m+n)^2+(m-n)^2\right).
\]
因为 \(m\neq n\)，所以 \((m-n)^2>0\)，从而
\[
S_m+S_n>\frac{9}{2}d^2k^2=\frac{9}{2}S_k.
\]
因此 \(c=\dfrac{9}{2}\) 时题中的不等式对所有符合条件的正整数都成立.

再证明不能取更大的值. 对任意正奇数 \(k\)，取
\[
m=\frac{3k-1}{2},\qquad n=\frac{3k+1}{2}.
\]
则 \(m,n\) 是不相等的正整数，且 \(m+n=3k\)，并且
\[
\frac{S_m+S_n}{S_k}
=\frac{m^2+n^2}{k^2}
=\frac{9}{2}+\frac{1}{2k^2}.
\]
当正奇数 \(k\) 充分大时，上式小于任意给定的 \(c>\dfrac{9}{2}\)，此时不等式 \(S_m+S_n>cS_k\) 不成立. 所以 \(c\) 的最大值为 \(\dfrac{9}{2}\).
\end{enumerate}
\end{solution}

\begin{problem}
设 \(f(x)\) 是定义在区间 \((1,+\infty)\) 上的函数，其导函数为 \(f'(x)\). 如果存在实数 \(a\) 和函数 \(h(x)\)，其中 \(h(x)\) 对任意的 \(x\in(1,+\infty)\) 都有 \(h(x)>0\)，使得 \(f'(x)=h(x)\left(x^2-ax+1\right)\)，则称函数 \(f(x)\) 具有性质 \(P(a)\).

\begin{enumerate}
\item 设函数 \(f(x)=\ln x+\frac{b+2}{x+1}\)（\(x>1\)），其中 \(b\) 为实数.
\begin{enumerate}
\item 求证：函数 \(f(x)\) 具有性质 \(P(b)\)；
\item 求函数 \(f(x)\) 的单调区间.
\end{enumerate}
\item 已知函数 \(g(x)\) 具有性质 \(P(2)\)，给定 \(x_1,x_2\in(1,+\infty)\)，\(x_1<x_2\)，设 \(m\) 为实数，\(\alpha=mx_1+(1-m)x_2\)，\(\beta=(1-m)x_1+mx_2\)，且 \(\alpha>1\)，\(\beta>1\)，若 \(|g(\alpha)-g(\beta)|<|g(x_1)-g(x_2)|\)，求 \(m\) 的取值范围.
\end{enumerate}
\end{problem}

\begin{solution}
\begin{enumerate}
\item
\begin{enumerate}
\item 对 \(f(x)\) 求导，得
\[
f'(x)=\frac{1}{x}-\frac{b+2}{(x+1)^2}
=\frac{x^2-bx+1}{x(x+1)^2}.
\]
当 \(x>1\) 时，令
\[
h(x)=\frac{1}{x(x+1)^2}>0,
\]
则 \(f'(x)=h(x)(x^2-bx+1)\)，所以函数 \(f(x)\) 具有性质 \(P(b)\).

\item 因为 \(x(x+1)^2>0\)，所以 \(f'(x)\) 的符号由 \(x^2-bx+1\) 决定.

当 \(b\leqslant2\) 时，
\[
x^2-bx+1=(x-1)^2+(2-b)x>0\qquad(x>1),
\]
所以 \(f(x)\) 在 \((1,+\infty)\) 上单调递增.

当 \(b>2\) 时，方程 \(x^2-bx+1=0\) 的两个根为
\[
x_1=\frac{b-\sqrt{b^2-4}}{2},\qquad x_2=\frac{b+\sqrt{b^2-4}}{2}.
\]
两根之积为 \(1\)，且 \(x_1<1<x_2\). 因而在定义域内，\(x^2-bx+1<0\) 的范围为 \(1<x<x_2\)，大于零的范围为 \(x>x_2\). 所以函数 \(f(x)\) 在
\[
\left(1,\frac{b+\sqrt{b^2-4}}{2}\right)
\]
上单调递减，在
\[
\left(\frac{b+\sqrt{b^2-4}}{2},+\infty\right)
\]
上单调递增.
\end{enumerate}

\item 因为函数 \(g(x)\) 具有性质 \(P(2)\)，所以
\[
g'(x)=h(x)(x^2-2x+1)=h(x)(x-1)^2>0\qquad(x>1).
\]
因此 \(g(x)\) 在 \((1,+\infty)\) 上严格递增.

记 \(d=x_2-x_1>0\)，则
\[
\alpha=x_2-md,\qquad \beta=x_1+md.
\]
当 \(0<m<1\) 时，\(x_1<\alpha<x_2\)，且 \(x_1<\beta<x_2\). 所以
\[
|g(\alpha)-g(\beta)|<g(x_2)-g(x_1)=|g(x_1)-g(x_2)|.
\]

当 \(m=0\) 或 \(m=1\) 时，\(\{\alpha,\beta\}=\{x_1,x_2\}\)，题中的不等式变为等式. 当 \(m<0\) 时，\(\alpha>x_2\)，\(\beta<x_1\)，由 \(g\) 的严格递增性，
\[
|g(\alpha)-g(\beta)|>g(x_2)-g(x_1).
\]
当 \(m>1\) 时，\(\alpha<x_1\)，\(\beta>x_2\)，同样有
\[
|g(\alpha)-g(\beta)|>g(x_2)-g(x_1).
\]
故 \(m\) 的取值范围为 \(0<m<1\).
\end{enumerate}
\end{solution}

\section{附加题：解答题}

\begin{problem}[A]
如图，\(AB\) 是 \(\odot O\) 的直径，\(D\) 为 \(\odot O\) 上一点，过点 \(D\) 作 \(\odot O\) 的切线交 \(AB\) 延长线于点 \(C\). 若 \(DA=DC\)，求证：\(AB=2BC\).

\bitmapfigure[max width=0.42\linewidth]{img/2010/jiangsu/q21_fig1.png}
\end{problem}

\begin{solution}
设 \(AB=2R\)，\(BC=x\)，其中 \(R>0\)，\(x>0\). 建立直角坐标系，取
\[
A(0,0),\qquad B(2R,0),\qquad C(2R+x,0),\qquad D(u,v).
\]
圆 \(\odot O\) 的方程为
\[
(X-R)^2+Y^2=R^2,
\]
所以
\[
u^2+v^2=2Ru.
\]
圆在点 \(D\) 处的切线方程为
\[
(u-R)(X-u)+v(Y-v)=0.
\]
因为 \(C\) 在该切线上，代入 \(C(2R+x,0)\)，得
\[
(u-R)(2R+x-u)-v^2=0.
\]
再利用 \(v^2=2Ru-u^2\)，化简得
\[
u(R+x)=R(2R+x).
\]

由点 \(C\) 的切割线定理，
\[
DC^2=CA\cdot CB=(2R+x)x.
\]
又因为 \(DA=DC\)，所以
\[
2Ru=DA^2=DC^2=x(2R+x).
\]
代入 \(u=\dfrac{R(2R+x)}{R+x}\)，并约去正数 \(2R+x\)，得
\[
\frac{2R^2}{R+x}=x,\qquad x^2+Rx-2R^2=0.
\]
于是 \((x-R)(x+2R)=0\). 因为 \(x>0\)，所以 \(x=R\)，从而
\[
AB=2R=2x=2BC.
\]
\end{solution}

\reuseproblemnumber
\begin{problem}[B]
在平面直角坐标系 \(xOy\) 中，\(A(0,0)\)，\(B(-2,0)\)，\(C(-2,1)\)，设 \(k\) 为非零实数. 矩阵 \(M=\begin{pmatrix}k&0\\0&1\end{pmatrix}\)，\(N=\begin{pmatrix}0&1\\1&0\end{pmatrix}\)，点 \(A\)，\(B\)，\(C\) 在矩阵 \(MN\) 对应的变换下得到的点分别为 \(A_1\)，\(B_1\)，\(C_1\)，\(\triangle A_1B_1C_1\) 的面积是 \(\triangle ABC\) 面积的 \(2\) 倍，求实数 \(k\) 的值.
\end{problem}

\begin{solution}
矩阵 \(MN\) 对应的线性变换将平面图形的面积扩大为原来的
\[
\left|\det(MN)\right|=\left|\det M\det N\right|=|k|
\]
倍. 又
\[
S_{\triangle ABC}=\frac12\times BC\times2=1.
\]
由题设 \(S_{\triangle A_1B_1C_1}=2S_{\triangle ABC}\)，得
\[
|k|=2.
\]
所以 \(k=2\) 或 \(k=-2\).
\end{solution}

\reuseproblemnumber
\begin{problem}[C]
在极坐标系中，圆 \(\rho=2\cos\theta\) 与直线 \(3\rho\cos\theta+4\rho\sin\theta+a=0\) 相切，求实数 \(a\) 的值.
\end{problem}

\begin{solution}
由极坐标与直角坐标的关系 \(x=\rho\cos\theta\)，\(y=\rho\sin\theta\)，圆的方程化为
\[
x^2+y^2=2x,
\]
即 \((x-1)^2+y^2=1\)，所以该圆的圆心为 \((1,0)\)，半径为 \(1\). 直线方程化为 \(3x+4y+a=0\).
相切条件是圆心到直线的距离等于半径，即
\[
\frac{|3+a|}{\sqrt{3^2+4^2}}=1.
\]
所以 \(|a+3|=5\)，解得
\[
a=2\quad\text{或}\quad a=-8.
\]
\end{solution}

\reuseproblemnumber
\begin{problem}[D]
设 \(a\)，\(b\) 是非负实数，求证：\(a^3+b^3\geqslant \sqrt{ab}\left(a^2+b^2\right)\).
\end{problem}

\begin{solution}
当 \(b=0\) 时，右边为 \(0\)，不等式显然成立. 以下设 \(b>0\)，令
\[
t=\sqrt{\frac{a}{b}}\geqslant0,\qquad a=t^2b.
\]
原不等式两边同除以正数 \(b^3\)，等价于
\[
t^6+1\geqslant t(t^4+1).
\]
又 \(t^6+1-t^5-t=(t-1)(t^5-1)=(t-1)^2\left(t^4+t^3+t^2+t+1\right)\geqslant0\).
所以原不等式成立.
\end{solution}

\section{解答题}

\begin{problem}
某工厂生产甲，乙两种产品. 甲产品的一等品率为 \(80\%\)，二等品率为 \(20\%\)；乙产品的一等品率为 \(90\%\)，二等品率为 \(10\%\). 生产 \(1\) 件甲产品，若是一等品则获得利润 \(4\) 万元，若是二等品则亏损 \(1\) 万元；生产 \(1\) 件乙产品，若是一等品则获得利润 \(6\) 万元，若是二等品则亏损 \(2\) 万元. 设生产各种产品相互独立.

\begin{enumerate}
\item 记 \(X\)（单位：万元）为生产 \(1\) 件甲产品和 \(1\) 件乙产品可获得的总利润，求 \(X\) 的分布列；
\item 求生产 \(4\) 件甲产品所获得的利润不少于 \(10\) 万元的概率.
\end{enumerate}
\end{problem}

\begin{solution}
设生产 \(1\) 件甲产品的利润为 \(U\)，生产 \(1\) 件乙产品的利润为 \(V\). 则
\[
P(U=4)=\frac45,\quad P(U=-1)=\frac15,
\]
\[
P(V=6)=\frac9{10},\quad P(V=-2)=\frac1{10}.
\]
又 \(U\)，\(V\) 相互独立，且 \(X=U+V\).

\begin{enumerate}
\item \(X\) 的可能取值及其概率如下：
\begin{center}
\begin{tabular}{c|cccc}
\(X\) & \(-3\) & \(2\) & \(5\) & \(10\) \\
\hline
\(P\) & \(\dfrac1{50}\) & \(\dfrac2{25}\) & \(\dfrac9{50}\) & \(\dfrac{18}{25}\)
\end{tabular}
\end{center}
其中
\[
P(X=-3)=P(U=-1,V=-2)=\frac15\times\frac1{10}=\frac1{50},
\]
\[
P(X=2)=P(U=4,V=-2)=\frac45\times\frac1{10}=\frac2{25},
\]
\[
P(X=5)=P(U=-1,V=6)=\frac15\times\frac9{10}=\frac9{50},
\]
\[
P(X=10)=P(U=4,V=6)=\frac45\times\frac9{10}=\frac{18}{25}.
\]
因此上表为 \(X\) 的分布列.

\item 设生产的 \(4\) 件甲产品中一等品的件数为 \(K\)，则 \(K\) 服从参数为 \(4\)，\(\dfrac45\) 的二项分布. 这 \(4\) 件甲产品的总利润为
\[
4K-(4-K)=5K-4.
\]
由 \(5K-4\geqslant10\)，得 \(K\geqslant3\). 因而所求概率为
\[
P(K\geqslant3)
=\mathrm C_4^3\left(\frac45\right)^3\frac15+\left(\frac45\right)^4
=\frac{512}{625}.
\]
\end{enumerate}
\end{solution}

\begin{problem}
已知 \(\triangle ABC\) 的三边长都是有理数.

\begin{enumerate}
\item 求证：\(\cos A\) 是有理数；
\item 求证：对任意正整数 \(n\)，\(\cos nA\) 是有理数.
\end{enumerate}
\end{problem}

\begin{solution}
\begin{enumerate}
\item 设三角形 \(ABC\) 的三边长分别为 \(a\)，\(b\)，\(c\). 由余弦定理，
\[
\cos A=\frac{b^2+c^2-a^2}{2bc}.
\]
因为 \(a\)，\(b\)，\(c\) 均为有理数，且 \(b,c>0\)，所以分子和分母都是有理数，且分母不为零. 有理数对加法、减法、乘法和除法（除数不为零）封闭，故 \(\cos A\) 是有理数.

\item 记 \(r=\cos A\)，则 \(r\) 是有理数. 当 \(n=1\) 时，\(\cos A=r\) 是有理数；当 \(n=2\) 时，
\[
\cos 2A=2\cos^2A-1=2r^2-1
\]
也是有理数.

对任意整数 \(k\geqslant2\)，由余弦和差公式可得
\[
2\cos A\cos kA=\cos(k+1)A+\cos(k-1)A.
\]
因此
\[
\cos(k+1)A=2\cos A\cos kA-\cos(k-1)A.
\]
假设 \(\cos(k-1)A\) 和 \(\cos kA\) 都是有理数，则由上式及 \(\cos A=r\in\mathbb{Q}\)，可知 \(\cos(k+1)A\) 也是有理数. 结合 \(n=1\) 和 \(n=2\) 时结论成立，利用数学归纳法可得：对任意正整数 \(n\)，\(\cos nA\) 都是有理数.
\end{enumerate}
\end{solution}
